\begin{textsample}{POS Dim 2 – ma – Score 29.00 – ma/ma\_em\_55.txt}  \label{ex:f2_pos_047}
5.3.
Formação inicial e continuada ( FIC ) ou \textbf{qualificação} profissional Os \textbf{cursos} de formação inicial e continuada ou de \textbf{qualificação} profissional para o trabalho consistem em processo de formação e desenvolvimento de competências e habilidades de um determinado perfil profissional \textbf{previsto} pelo mercado de trabalho ( BRASIL, 2018 ).
Podem ser organizados com \textbf{carga} \textbf{horária} mínima de 160 horas, considerando o \textbf{previsto} em \textbf{legislação} e as \textbf{cargas} \textbf{horárias} das \textbf{habilitações} profissionais especificadas no CNTC.
Também podem constituir \textbf{itinerários} formativos integrantes de uma matriz curricular de determinado \textbf{curso} \textbf{técnico} de nível médio, como possibilidade de saída intermediária.
De forma que, após conclusão do \textbf{curso} FIC, o estudante é certificado de acordo com normativas nacionais e estaduais.
Poderão, ainda, assumir formas de \textbf{cursos} especiais abertos à comunidade em geral, condicionando a matrícula à capacidade de aproveitamento dos estudantes.
I - Programa de aprendizagem Conforme Resolução nº 3/2018, programa de aprendizagem consiste em organizações curriculares articuladas de \textbf{cursos} de educação profissional e \textbf{técnica}, que viabilizem aproveitamentos curriculares em uma determinada trilha formativa profissional.
A \textbf{oferta} de programa de aprendizagem pode estar vinculada tanto à \textbf{habilitação} profissional \textbf{técnica} de nível médio quanto à \textbf{qualificação} profissional ( BRASIL, 2018 ).
Essa forma de \textbf{oferta} objetiva apoiar percursos formativos que corroborem para a formação integral dos estudantes, incorrendo em sua inserção prévia no mundo trabalho, por meio de vivências ou práticas profissionais em empresas durante o período de formação escolar.
Regulamentado pela Lei nº 10.097/2000, que determina que as empresas de médio e grande porte contratem um número de aprendizes ( entre 14 e 24 anos incompletos ) equivalente a um mínimo de 5 \% e um máximo de 15 \% do total de empregados, em funções que demandem \textbf{qualificação}, o programa de aprendizagem surge como oportunidade de formação, inserção no mundo produtivo, realização de projetos de vidas, bem como consolidação de valores.
Nesse sentido, possui como finalidade precípua o combate à evasão escolar ; a \textbf{garantia} de que \textbf{adolescentes} e jovens apliquem em atividades práticas os conhecimentos adquiridos no processo formativo ; \textbf{qualificação} para a atuação profissional ; o aumento da empregabilidade ; o diálogo social e ainda a conciliação entre estudos, trabalho e vida familiar.
Modalidade de \textbf{oferta} Segundo a Resolução nº 3/2018, podem ser consideradas como \textbf{carga} \textbf{horária} do ensino médio as atividades com fins pedagógicos e critérios estabelecidos pelos sistemas de ensino, tais como : “ (... ) aulas, \textbf{cursos}, estágios, oficinas, trabalho supervisionado, atividades de extensão, pesquisa de campo, iniciação científica, aprendizagem profissional, participação em trabalhos voluntários ”, entre outras ( BRASIL, 2018, p. 11 ), podendo, ainda, ser computadas como \textbf{certificações} complementares, constando no histórico escolar.
Essas atividades podem ocorrer de forma presencial, auxiliada ou não por ferramentas tecnológicas, ou a distância.
Na forma a distância, as atividades podem contemplar até 20 \% ( vinte por cento ) da \textbf{carga} \textbf{horária} total, incorrendo prioritariamente nos \textbf{itinerários} formativos do \textbf{currículo}, mediante suporte tecnológico e pedagógico adequados.
A depender dos sistemas de ensino, essa forma poderá se expandir para até 30 \% ( trinta por cento ), no ensino médio noturno, conforme Resolução CNE/CEB nº 3/2018 e Resolução CEE/MA nº 277/2020.
Estrutura dos \textbf{itinerários} de educação \textbf{técnica} e profissional da rede No que se refere à formação \textbf{técnica} e profissional, a organização dos \textbf{itinerários} ocorre por meio da integração dos eixos estruturantes ( processos criativos, investigação científica, mediação e intervenção sociocultural e empreendedorismo ), almejando ampliar habilidades pertinentes ao pensar e fazer científico, à prática da convivência e atuação em sociedade, além da \textbf{gestão} de ações empreendedoras em torno da concretização do projeto de vida dos estudantes.
Nessa direção, apresenta -se como possibilidade de percursos formativos na educação profissional e \textbf{técnica} na rede pública de ensino do Maranhão dispostas no organograma que segue.
Essa organização curricular, além dos componentes que integram à estruturação do aprofundamento de formação \textbf{técnica} e profissional em áreas específicas, também compreende : Figura 7 – Percursos formativos na educação profissional e \textbf{técnica} na rede pública de ensino do Maranhão - Ensino médio articulado - Formação geral básica - Itinerário \textbf{técnico} e profissional - Integral integrado - Modalidades integradas - Parcial integrado concomitante - \textbf{Habilitação} profissional \textbf{técnica} de nível médio ( 800, 1.000, 1.200 horas ) - \textbf{Qualificação} profissional \textbf{técnica} de nível médio ( 800, 1.000, 1.200 horas ) - Programa de aprendizagem - Conjunto de FICs articuladas - \textbf{Habilitação} profissional \textbf{técnica} de nível médio ( 800, 1.000, 1.200 horas ) - \textbf{Habilitação} profissional \textbf{técnica} de nível médio ( 800, 1.000, 1.200 horas ) - \textbf{Qualificação} profissional \textbf{técnica} de nível médio ( 800, 1.000, 1.200 horas ) - Conjunto de FICs articuladas - Formação diversificada/preparação básica para o trabalho - Formação profissional Fonte : Organizado pelos autores, 2020 Figura 8 – Estruturação do aprofundamento de formação \textbf{técnica} e profissional em áreas específicas - Escolas em tempo integral - \textbf{Eletivas} + Projeto de Vida + Pós-médio + Estudo Orientado e Avaliação Semanal + Tutoria - Escolas em tempo parcial - Parte diversificada - \textbf{Eletivas} + Projeto de Vida + Pós-médio + Tutoria - Projetos empreendedores - Modalidades integradas - Projeto de Vida - Projetos empreendedores - Práticas experimentais + Projetos empreendedores + Projetos de \textbf{Corresponsabilidade} Social - Formação básica para o trabalho - Projetos de \textbf{Corresponsabilidade} Social + Empreendedorismo + Saúde e Segurança no Trabalho - Empreendedorismo + Pesquisa Científica + Linguagem, Trabalho e Tecnologia + Intervenção Sociocultural - Vivência Profissional ou Estágio ou TCC - Formação específica do \textbf{curso} - Vivência Profissional ou Estágio ou TCC - Vivência Profissional ou Estágio ou TCC Fonte : Organizado pelos autores, 2020 O projeto de vida, enquanto componente curricular, integrará todos os \textbf{itinerários} formativos do \textbf{currículo} do ensino médio do estado do Maranhão, de modo a ser estruturado e trabalhado ao longo dessa etapa de ensino, de forma transversal ao cotidiano escolar ( BRASIL, 2019 ).
A metodologia central do projeto de vida será a de aprendizagem por projetos, personalizando e valorizando as vivências dos estudantes, visando ao desenvolvimento de competências e habilidades para a vida contemporânea, com base em valores e responsabilidade social e ambiental.
No que se refere ao \textbf{itinerário} \textbf{técnico} e profissional, o componente curricular projeto de vida se integrará à \textbf{carga} \textbf{horária} das diversas formas de \textbf{ofertas}.
Formação no trabalho Para cumprir exigências curriculares do ensino médio, os sistemas de ensino poderão estabelecer critérios próprios e específicos para reconhecer competências dos estudantes tanto da formação geral básica quanto dos \textbf{itinerários} formativos do \textbf{currículo} ( BRASIL, 2018 ).
Nesse sentido, para efetivação da \textbf{certificação} profissional, que consiste no processo de avaliação, reconhecimento e \textbf{certificação} de saberes adquiridos na educação profissional, inclusive no trabalho, para fins de \textbf{prosseguimento} ou conclusão de estudos, nos termos do \textbf{art}. 41 da LDB nº 9.394/96, deverão ser adotados critérios específicos validados pelo Conselho Estadual do Educação do Maranhão ( CEE -MA ).
No que se refere ao reconhecimento de \textbf{notório} saber, a Lei nº 13.415/2017 estabelece que os profissionais com \textbf{notório} saber, uma vez reconhecidos pelos respectivos sistemas de ensino, estarão habilitados a ministrar objetos do conhecimento de áreas correlatas à sua formação ou experiência profissional atestada.
O estágio supervisionado, que consiste em atividade pedagógica, é regulamentado nacionalmente pela Lei nº 11.788, de 25 de setembro de 2008.
Sua concretização congrega conhecimentos acadêmicos com experiências vivenciadas no ambiente de trabalho em áreas comuns às da formação \textbf{técnica} e profissional do estudante, favorecendo, com isso, o aprofundamento dos conhecimentos por meio da prática profissional.
Constitui etapa de fundamental importância para o processo de desenvolvimento e aprendizagem do estudante que almeja a concretização de seus objetivos por meio de uma \textbf{habilitação} ou \textbf{qualificação} \textbf{técnica}.
O estágio caracteriza -se como uma atividade específica de vivência profissional, que mantém relação estreita com a transição pela qual o estudante passa de um ambiente escolar para o ambiente no qual exercerá sua futura profissão.
Nesse sentido, o estágio supervisionado tem por objetivo a aplicação de conhecimentos teóricos e práticos adquiridos pelo estudante no decorrer do \textbf{curso}.
O estágio curricular ocorrerá em \textbf{conformidade} com o plano de \textbf{curso} de cada \textbf{curso} de \textbf{habilitação} \textbf{técnica} e profissional.
Essa ação pedagógica é regulamentada pela Lei nº 11.788, de 25 de setembro de 2008, pela Lei de \textbf{Diretrizes} e Bases da Educação Nacional nº 9.394/96, pela Resolução CNE/CP nº 1, de 5 de \textbf{janeiro} de 2021, que define as \textbf{Diretrizes} Curriculares Nacionais Gerais para a Educação Profissional e Tecnológica, além da Resolução nº 120/2013-CEE/MA, que estabelece normas para a educação profissional \textbf{técnica} de nível médio no sistema estadual de ensino do Maranhão.

% matched lemmas: adolescente, art, carga, certificação, conformidade, corresponsabilidade, currículo, curso, diretriz, eletivo, garantia, gestão, habilitação, horário, itinerário, janeiro, legislação, notório, oferta, prever, previsto, prosseguimento, qualificação, técnico
\end{textsample}
